    \begin{resumo}
        O objetivo deste trabalho é analisar a interação prática entre a Teoria
        da Probabilidade, o Teorema de Bayes, o Valor Esperado e a Variância na
        tomada de decisões sob incerteza. Para isso, foi realizada uma revisão
        teórica bibliográfica aliada a uma simulação computacional prática para
        modelar o comportamento e o risco de eventos discretos. Os resultados
        demonstram que esses quatro conceitos estatísticos não operam
        isoladamente, comprovando que formam um fluxo sequencial e dependente
        para a análise segura e assertiva de cenários probabilísticos.
    \end{resumo}
